% Part: first-order-logic
% Chapter: syntax-and-semantics
% Section: unique-readability

\documentclass[../../../include/open-logic-section]{subfiles}

\begin{document}

\olfileid{fol}{syn}{unq}

\olsection{ఏకార్థ పఠనీయత}

\begin{explain}
మనం \tetoken{సూత్రాలను}{!!{formula}s} నిర్వచించిన విధానం ప్రతి \tetoken{సూత్రానికి}{!!{formula}} ఒక
\emph{ఏకార్థ పఠనం} ఉంటుందని హామీ ఇస్తుంది. అంటే, \tetoken{సూత్రాలకు}{!!{formula}s} మన
నిర్మాణ నియమాల ప్రకారం దాన్ని నిర్మించడానికి సారాంశంగా ఒకే మార్గమూ,
దానికి ``అర్థం చెప్పడానికి'' ఒకే మార్గమూ ఉంటాయి. అలా లేకపోతే మనకు
సందిగ్ధమైన \tetoken{సూత్రాలు}{!!{formula}s}, అంటే ఒకటికంటే ఎక్కువ పఠనాలు లేదా అర్థాలు గల
\tetoken{సూత్రాలు}{!!{formula}s} ఉండేవి---దాన్ని మనం నివారించాలనుకోవడం స్పష్టమే. మరింత
ముఖ్యంగా, ఈ ధర్మం లేకపోతే మనం ఇవ్వబోయే నిర్వచనాలు, నిరూపణల్లో అధిక
భాగం పనిచేయదు.

ఏకార్థ పఠనీయతకు హామీ ఇవ్వని చెడ్డ నిర్మాణ నియమాలను \tetoken{సూత్రాలకు}{!!{formula}s}
ఇచ్చి ఉంటే ఏమయ్యేదో చూడడం ద్వారా దీన్ని బహుశా అత్యంత స్పష్టంగా
చూడవచ్చు. ఉదాహరణకు, సంయోజకాల నిర్మాణ నియమాల్లో కుండలీకరణ చిహ్నాలను
మర్చిపోయి, కిందిదాన్ని అనుమతించి ఉండవచ్చు:
\begin{quote}
$!A$, $!B$లు \tetoken{సూత్రాలు}{!!{formula}s} అయితే, $!A \lif !B$ కూడా ఒక సూత్రమే.
\end{quote}
పరమాణు సూత్రం~$!D$తో మొదలుపెడితే, దీని వల్ల $!D \lif !D$ను
ఏర్పరచవచ్చు. దీని నుంచీ $!D$ నుంచీ $!D \lif !D \lif !D$ను పొందవచ్చు.
కానీ దీన్ని చేయడానికి రెండు మార్గాలు ఉన్నాయి:
\begin{enumerate}
\item $!D$ను $!A$గానూ, $!D \lif !D$ను $!B$గానూ తీసుకుంటాం.
\item $!A$ను $!D \lif !D$గానూ, $!B$ను~$!D$గానూ తీసుకుంటాం.
\end{enumerate}
దానికి అనుగుణంగా \tetoken{సూత్రం}{!!{formula}}~$!D \lif !D \lif !D$ను ``చదవడానికి''
రెండు మార్గాలు ఉన్నాయి. అది $!B \lif !C$ రూపంలో ఉంటుంది; ఇందులో
$!B$ అనేది $!D$, $!C$ అనేది $!D \lif !D$. కానీ అది $!B \lif !C$
రూపంలో \emph{కూడా} ఉంటుంది; అప్పుడు $!B$ అనేది $!D \lif !D$, $!C$
అనేది~$!D$.

ఇలా జరిగితే మన నిర్వచనాలు ఎల్లప్పుడూ పనిచేయవు. ఉదాహరణకు, ఒక సూత్రపు
\tetoken{ప్రధాన సంయోజకాన్ని}{!!{main operator}} నిర్వచించేటప్పుడు ఇలా అంటాం: $!B \lif !C$ రూపపు
సూత్రంలో సూచించిన~$\lif$ ఘటనే \tetoken{ప్రధాన సంయోజకం}{!!{main operator}}. కానీ
$!D \lif !D \lif !D$ సూత్రాన్ని పైన చెప్పిన రెండు వేర్వేరు మార్గాల్లో
$!B \lif !C$తో సరిపోల్చగలిగితే, ఒక సందర్భంలో మొదటి $\lif$ ఘటనే
\tetoken{ప్రధాన సంయోజకం}{!!{main operator}} అవుతుంది; రెండవ సందర్భంలో రెండవ ఘటనే అవుతుంది.
అయితే \tetoken{ప్రధాన సంయోజకం}{!!{main
  operator}} అనేది \tetoken{సూత్రంపై}{!!{formula}} ఒక \emph{ప్రమేయం} కావాలని
మన ఉద్దేశం. అంటే ప్రతి \tetoken{సూత్రానికి}{!!{formula}} సరిగ్గా ఒక్క \tetoken{ప్రధాన సంయోజకం}{!!{main operator}}
ఘటన ఉండాలి.
\end{explain}

\begin{lem}
\tetoken{ఒక సూత్రం}{!!a{formula}}~$!A$లోని ఎడమ, కుడి కుండలీకరణ చిహ్నాల సంఖ్యలు సమానం.
\end{lem}

\begin{proof}
$!A$ను నిర్మించిన విధానంపై ఆగమనం ద్వారా దీన్ని నిరూపిస్తాం. దీనికి
రెండు విషయాలు కావాలి: (a) ముందుగా అన్ని పరమాణు సూత్రాలకూ పరిశీలిస్తున్న
ధర్మం ఉందని నిరూపించాలి (ఆగమన ఆధారం). (b) తరువాత, ఇచ్చిన సూత్రాల నుంచి
కొత్త సూత్రాలను నిర్మించినప్పుడు, పాత వాటికి ఆ ధర్మం ఉంటే కొత్త వాటికీ
ఉంటుందని నిరూపించాలి.

$l(!A)$ అనేది $!A$లోని ఎడమ కుండలీకరణ చిహ్నాల సంఖ్య, $r(!A)$ అనేది
కుడి కుండలీకరణ చిహ్నాల సంఖ్య అనుకుందాం. అదే విధంగా $l(t)$, $r(t)$లు
పదం~$t$లోని ఎడమ, కుడి కుండలీకరణ చిహ్నాల సంఖ్యలు.

\begin{prob}
    ఏ పదం~$t$కైనా $l(t) = r(t)$ అని నిరూపించండి.
\end{prob}

\begin{enumerate}
\tagitem{prvFalse}{\indcase{!A}{\lfalse}{$\indfrm$లో $0$ ఎడమ,
    $0$ కుడి కుండలీకరణ చిహ్నాలు ఉంటాయి.}}{}

\tagitem{prvTrue}{\indcase{!A}{\ltrue}{$\indfrm$లో $0$ ఎడమ,
    $0$ కుడి కుండలీకరణ చిహ్నాలు ఉంటాయి.}}{}

\item \indcase{!A}{\Atom{R}{t_1,\dots,t_n}}{$l(\indfrm) = 1 + l(t_1) +
  \dots + l(t_n) = 1 + r(t_1) + \dots + r(t_n) = r(\indfrm)$.
  ఇక్కడ ఏ పదం~$t$కైనా $l(t) = r(t)$ అనే, సాధనగా విడిచిపెట్టిన
  వాస్తవాన్ని ఉపయోగించాం.}

\item \indcase{!A}{\eq[t_1][t_2]}{$l(\indfrm) = l(t_1) + l(t_2) =
  r(t_1) + r(t_2) = r(\indfrm)$.}

\tagitem{prvNot}{\indcase{!A}{\lnot !B}{ఆగమన పరికల్పన ప్రకారం,
    $l(!B) = r(!B)$. అందువల్ల $l(\indfrm) = l(!B) = r(!B) =
    r(\indfrm)$.}}{}

\item \indcase{!A}{(!B \ast !C)}{ఆగమన పరికల్పన ప్రకారం $l(!B) =
  r(!B)$, $l(!C) = r(!C)$. అందువల్ల $l(\indfrm) = 1 + l(!B) + l(!C) =
  1 + r(!B) + r(!C) = r(\indfrm)$.}

\tagitem{prvAll}{\indcase{!A}{\lforall[x][!B]}{ఆగమన పరికల్పన
    ప్రకారం $l(!B) = r(!B)$. అందువల్ల $l(\indfrm) = l(!B) = r(!B) =
    r(\indfrm)$.}}{}

% Print case for \lexists if prvEx and not prvAll
\tagitem{prvAll,notprvEx}{}{\indcase{!A}{\lexists[x][!B]}{ఆగమన
    పరికల్పన ప్రకారం $l(!B) = r(!B)$. అందువల్ల
    $l(\indfrm) = l(!B) = r(!B) = r(\indfrm)$.}}

% Just say similarly if prvEx and prvAll
\tagitem{notprvAll,notprvEx}{}{\indcase{!A}{\lexists[x][!B]}{అదే విధంగా.}}
\end{enumerate}
\end{proof}

\begin{defn}[నిజ ఆద్యఖండం]
సంకేతమాల~$!B$కు శూన్యేతర సంకేతమాలను సంయోజించినప్పుడు సంకేతమాల~$!A$
వస్తే, $!B$ను $!A$ యొక్క \emph{నిజ ఆద్యఖండం} అంటాం.
\end{defn}

\begin{lem}\ollabel{lem:no-prefix}
$!A$ \tetoken{ఒక సూత్రం}{!!a{formula}}, $!B$ అనేది $!A$కు నిజ ఆద్యఖండం అయితే, $!B$
\tetoken{ఒక సూత్రం}{!!a{formula}} కాదు.
\end{lem}

\begin{proof}
సాధన.
\end{proof}

\begin{prob}
\olref[fol][syn][unq]{lem:no-prefix}ను నిరూపించండి.
\end{prob}

\begin{prop}
\ollabel{prop:unique-atomic}
$!A$ పరమాణు \tetoken{సూత్రం}{!!{formula}} అయితే, కింది షరతుల్లో ఒక్కటి, ఒక్కటి మాత్రమే
దానికి వర్తిస్తుంది.
\begin{enumerate}
\tagitem{prvFalse}{$!A \ident \lfalse$.}{}
\tagitem{prvTrue}{$!A \ident \ltrue$.}{}
\item $!A \ident \Atom{R}{t_1,\dots,t_n}$; ఇక్కడ $R$ ఒక $n$-స్థాన
  \tetoken{విధేయ సంకేతం}{!!{predicate}}, $t_1$, \dots, $t_n$లు పదాలు, ఇంకా $R$,
  $t_1$, \dots, $t_n$లలో ప్రతిదీ ఏకైకంగా నిర్ణయితమవుతుంది.
\item $!A \ident \eq[t_1][t_2]$; ఇక్కడ $t_1$, $t_2$లు ఏకైకంగా
  నిర్ణయితమైన పదాలు.
\end{enumerate}
\end{prop}

\begin{proof}
సాధన.
\end{proof}

\begin{prob}
\olref[fol][syn][unq]{prop:unique-atomic}ను నిరూపించండి (సూచన:
పదాల కోసం \olref[fol][syn][unq]{lem:no-prefix}కు ఒక రూపాన్ని
ప్రతిపాదించి నిరూపించండి).
\end{prob}

\begin{prop}[ఏకార్థ పఠనీయత]
ప్రతి \tetoken{సూత్రం}{!!{formula}} కింది షరతుల్లో ఒక్కదాన్ని, ఒక్కదాన్ని మాత్రమే తీరుస్తుంది.
\begin{enumerate}
\item $!A$ పరమాణువు.

\tagitem{prvNot}{$!A$, $\lnot !B$ రూపంలో ఉంటుంది.}{}

\tagitem{prvAnd}{$!A$, $(!B \land !C)$ రూపంలో ఉంటుంది.}{}

\tagitem{prvOr}{$!A$, $(!B \lor !C)$ రూపంలో ఉంటుంది.}{}

\tagitem{prvIf}{$!A$, $(!B \lif !C)$ రూపంలో ఉంటుంది.}{}

\tagitem{prvIff}{$!A$, $(!B \liff !C)$ రూపంలో ఉంటుంది.}{}

\tagitem{prvAll}{$!A$, $\lforall[x][!B]$ రూపంలో ఉంటుంది.}{}

\tagitem{prvEx}{$!A$, $\lexists[x][!B]$ రూపంలో ఉంటుంది.}{}
\end{enumerate}
అంతేకాక, ప్రతి సందర్భంలో $!B$, లేదా $!B$, $!C$లు ఏకైకంగా
నిర్ణయితమవుతాయి. ఉదాహరణకు, $!A$ కింది రెండు రూపాల్లోనూ ఉండేలా
వేర్వేరు జతలు $!B$, $!C$ మరియు $!B'$, $!C'$ ఉండవని దీని అర్థం:
\iftag{prvIf}{$(!B \lif !C)$ మరియు $(!B' \lif !C')$.}{
    \iftag{prvOr}{$(!B \lor !C)$ మరియు $(!B' \lor !C')$.}{
      \iftag{prvAnd}{$(!B \land !C)$ మరియు $(!B' \land !C')$.}{}}}
\end{prop}

\begin{proof}
\tetoken{ఒక సూత్రం}{!!a{formula}} పరమాణువు కాకపోతే, అది తప్పనిసరిగా ఎడమ కుండలీకరణ
చిహ్నం~(తో, \iftag{prvNot}{$\lnot$తో,}{} లేదా ఒక పరిమాణీకరణికంతో
మొదలవ్వాలని నిర్మాణ నియమాలు కోరుతాయి. మరోవైపు, కింది సంకేతాల్లో
దేనితోనైనా మొదలయ్యే ప్రతి \tetoken{సూత్రం}{!!{formula}} తప్పనిసరిగా పరమాణువు:
\tetoken{ఒక విధేయ సంకేతం}{!!a{predicate}}, \tetoken{ఒక ప్రమేయ సంకేతం}{!!a{function}}, \tetoken{ఒక స్థిరసంకేతం}{!!a{constant}}\iftag{prvFalse}{,
$\lfalse$}{}\iftag{prvTrue}{, $\ltrue$}{}.

కాబట్టి $!A$ ఒకవైపు $(!B \ast !C)$ రూపంలోనూ, మరోవైపు
$(!B' \mathbin{\ast'} !C')$ రూపంలోనూ ఉంటే, $!B \ident !B'$,
$!C \ident !C'$, $\ast = {\ast'}$ అని మాత్రమే చూపాలి.

$!A \ident (!B \ast !C)$, $!A \ident (!B'
\mathbin{\ast'} !C')$ రెండూ అనుకుందాం. అప్పుడు $!B \ident !B'$
లేదా అలా కాదు. అలా అయితే, అవి $!A$లో ఒకే చోట మొదలై సమాన పొడవు గల
ఉపసంకేతమాలలు కనుక $\ast = {\ast'}$, $!C \ident !C'$ అని స్పష్టం.
మిగిలిన సందర్భం $!B \not\ident !B'$. $!B$, $!B'$ రెండూ $!A$లో
ఒకే చోట మొదలయ్యే ఉపసంకేతమాలలు కనుక, వాటిలో ఒకటి మరొకదానికి నిజ
ఆద్యఖండం అయి ఉండాలి. కానీ \olref{lem:no-prefix} ప్రకారం ఇది అసాధ్యం.
\end{proof}


\end{document}
